\chapter{§1.2 四则运算}
\label{sec:1.2}


\prerequisite{§1.1 整数的认识}

\gradelevel{1-4 年级}


\section*{加法与减法}


\begin{explore}


书包里有 $5$ 支铅笔，妈妈又给了你 $3$ 支。现在一共有多少支？如果你用掉了 $2$ 支，还剩多少支？

\end{explore}

\begin{understand}


\textbf{加法}表示把两个数合在一起。 $5 + 3 = 8$ ，其中 $5$ 和 $3$ 叫做\textbf{加数}， $8$ 叫做\textbf{和}。

\textbf{减法}是加法的逆运算，表示从一个数里去掉一部分。 $8 - 2 = 6$ ，其中 $8$ 叫做\textbf{被减数}， $2$ 叫做\textbf{减数}， $6$ 叫做\textbf{差}。

\textbf{竖式计算}：相同数位对齐，从个位算起。
\begin{itemize}
  \item 加法：满十向前一位进一。
  \item 减法：不够减时向前一位借一当十。
\end{itemize}


加法竖式示例（ $47 + 35$ ）：

\[
\begin{array}{rl}
&\phantom{+}47 \\
+\enspace&\phantom{0}35 \\
\hline
&\phantom{+}82
\end{array}
\]


个位： $7 + 5 = 12$ ，写 $2$ 进 $1$ 。十位： $4 + 3 + 1 = 8$ 。

减法竖式示例（ $82 - 35$ ）：

\[
\begin{array}{rl}
&\phantom{-}82 \\
-\enspace&\phantom{0}35 \\
\hline
&\phantom{-}47
\end{array}
\]


个位： $2 < 5$ ，向十位借一， $12 - 5 = 7$ 。十位： $7 - 3 = 4$ （借走一个后 $8$ 变 $7$ ）。

\end{understand}

\begin{workedexamples}


\textbf{例 1.} 计算 $256 + 178$ 。

\textbf{解：}

\[
\begin{array}{rl}
&\phantom{+}256 \\
+\enspace&\phantom{0}178 \\
\hline
&\phantom{+}434
\end{array}
\]


\begin{itemize}
  \item 个位： $6 + 8 = 14$ ，写 $4$ 进 $1$ 。
  \item 十位： $5 + 7 + 1 = 13$ ，写 $3$ 进 $1$ 。
  \item 百位： $2 + 1 + 1 = 4$ 。
\end{itemize}


所以 $256 + 178 = 434$ 。

\textbf{例 2.} 计算 $503 - 267$ 。

\textbf{解：}

\[
\begin{array}{rl}
&\phantom{-}503 \\
-\enspace&\phantom{0}267 \\
\hline
&\phantom{-}236
\end{array}
\]


\begin{itemize}
  \item 个位： $3 < 7$ ，向十位借一，但十位是 $0$ ，需要从百位借。百位借 $1$ 给十位得 $10$ ，十位再借 $1$ 给个位： $13 - 7 = 6$ 。
  \item 十位：借走后剩 $9$ ， $9 - 6 = 3$ 。
  \item 百位：借走后剩 $4$ ， $4 - 2 = 2$ 。
\end{itemize}


所以 $503 - 267 = 236$ 。

\end{workedexamples}

\begin{keytakeaway}


\begin{itemize}
  \item 加法和减法互为逆运算。
  \item 竖式计算时相同数位对齐，从个位算起。
  \item 加法满十进一，减法不够减时借一当十。
\end{itemize}


\end{keytakeaway}

\begin{practice}


\begin{enumerate}
  \item 计算： $349 + 256$ 。
  \item 计算： $800 - 453$ 。
  \item 学校合唱团有 $128$ 名同学，又加入了 $37$ 名，后来有 $15$ 名同学退出了。现在合唱团有多少人？
\end{enumerate}


\end{practice}

\section*{乘法与除法}


\begin{explore}


每个小组有 $4$ 个人，教室里有 $6$ 个小组，一共有多少个人？你可以算 $4 + 4 + 4 + 4 + 4 + 4$ ，但有没有更快的方法？

\end{explore}

\begin{understand}


\textbf{乘法}是求几个相同加数之和的简便运算。 $4 \times 6 = 24$ ，其中 $4$ 和 $6$ 叫做\textbf{因数}， $24$ 叫做\textbf{积}。

\textbf{除法}是乘法的逆运算，表示平均分。 $24 \div 6 = 4$ ，其中 $24$ 叫做\textbf{被除数}， $6$ 叫做\textbf{除数}， $4$ 叫做\textbf{商}。

\end{understand}

\begin{quote}
注意：除数不能为 $0$ 。
\end{quote}


\textbf{乘法竖式}示例（ $23 \times 14$ ）：

\[
\begin{array}{rl}
&\phantom{\times}23 \\
\times\enspace&\phantom{0}14 \\
\hline
&\phantom{\times}92 \quad \leftarrow 23 \times 4 \\
&\phantom{0}23\phantom{0} \quad \leftarrow 23 \times 10 \\
\hline
&\phantom{0}322
\end{array}
\]


\textbf{除法竖式}示例（ $96 \div 4 = 24$ ）：

先看十位： $9 \div 4 = 2$ 余 $1$ ，商的十位写 $2$ ，余下的 $1$ 和个位的 $6$ 合成 $16$ 。再算 $16 \div 4 = 4$ ，商的个位写 $4$ 。

\begin{workedexamples}


\textbf{例 1.} 计算 $36 \times 25$ 。

\textbf{解：}

\[
\begin{array}{rl}
&\phantom{\times}36 \\
\times\enspace&\phantom{0}25 \\
\hline
&\phantom{0}180 \quad \leftarrow 36 \times 5 \\
&\phantom{0}72\phantom{0} \quad \leftarrow 36 \times 20 \\
\hline
&\phantom{0}900
\end{array}
\]


所以 $36 \times 25 = 900$ 。

\textbf{例 2.} 计算 $756 \div 6$ 。

\textbf{解：}
\begin{itemize}
  \item $7 \div 6 = 1$ 余 $1$ ，商的百位写 $1$ 。
  \item 余 $1$ 和十位的 $5$ 合成 $15$ ， $15 \div 6 = 2$ 余 $3$ ，商的十位写 $2$ 。
  \item 余 $3$ 和个位的 $6$ 合成 $36$ ， $36 \div 6 = 6$ ，商的个位写 $6$ 。
\end{itemize}


所以 $756 \div 6 = 126$ 。

\end{workedexamples}

\begin{keytakeaway}


\begin{itemize}
  \item 乘法是相同加数求和的简便运算，除法是乘法的逆运算。
  \item 乘法竖式按位相乘再相加。
  \item 除法从最高位除起，余数要带到下一位继续除。
\end{itemize}


\end{keytakeaway}

\begin{practice}


\begin{enumerate}
  \item 计算： $48 \times 37$ 。
  \item 计算： $864 \div 4$ 。
  \item 一盒彩笔有 $12$ 支，老师买了 $15$ 盒，一共有多少支彩笔？
  \item 学校有 $312$ 本课外书，平均分给 $8$ 个班级，每班分到多少本，还剩多少本？
\end{enumerate}


\end{practice}

\section*{混合运算与运算顺序}


\begin{explore}


小明去文具店买了 $3$ 本练习本，每本 $2$ 元，又买了一块 $5$ 元的橡皮。他一共花了多少钱？算式写成 $3 \times 2 + 5$ ，该先算哪一步？

\end{explore}

\begin{understand}


当一个算式里有多种运算时，需要约定\textbf{运算顺序}：

\begin{enumerate}
  \item \textbf{先乘除，后加减}。
  \item \textbf{有括号的先算括号里面的}。
  \item \textbf{同级运算从左到右}依次计算。
\end{enumerate}


例如：

\[
3 \times 2 + 5 = 6 + 5 = 11
\]


\[
3 \times (2 + 5) = 3 \times 7 = 21
\]


括号改变了运算顺序，结果完全不同。

\end{understand}

\begin{workedexamples}


\textbf{例 1.} 计算 $120 - 18 \times 5$ 。

\textbf{解：} 先算乘法，再算减法。

\[
120 - 18 \times 5 = 120 - 90 = 30
\]


\textbf{例 2.} 计算 $(120 - 18) \times 5$ 。

\textbf{解：} 先算括号里面的，再算乘法。

\[
(120 - 18) \times 5 = 102 \times 5 = 510
\]


\end{workedexamples}

\begin{keytakeaway}


\begin{itemize}
  \item 先乘除，后加减。
  \item 有括号先算括号里面的。
  \item 同级运算从左到右。
  \item 括号可以改变运算顺序。
\end{itemize}


\end{keytakeaway}

\begin{practice}


\begin{enumerate}
  \item 计算： $24 + 6 \times 8$ 。
  \item 计算： $(24 + 6) \times 8$ 。
  \item 计算： $100 - (25 + 15) \times 2$ 。
  \item 小红有 $50$ 元，买了 $4$ 支笔，每支 $6$ 元，还剩多少钱？写出综合算式并计算。
\end{enumerate}


\end{practice}



\begin{answer}


\textbf{练一练 加法与减法}

\begin{enumerate}
  \item $349 + 256 = 605$ 。
  \item $800 - 453 = 347$ 。
  \item $128 + 37 - 15 = 150$ （人）。
\end{enumerate}


\textbf{练一练 乘法与除法}

\begin{enumerate}
  \item $48 \times 37 = 1776$ 。
  \item $864 \div 4 = 216$ 。
  \item $12 \times 15 = 180$ （支）。
  \item $312 \div 8 = 39$ （本），每班分到 $39$ 本，没有剩余。
\end{enumerate}


\textbf{练一练 混合运算与运算顺序}

\begin{enumerate}
  \item $24 + 6 \times 8 = 24 + 48 = 72$ 。
  \item $(24 + 6) \times 8 = 30 \times 8 = 240$ 。
  \item $100 - (25 + 15) \times 2 = 100 - 40 \times 2 = 100 - 80 = 20$ 。
  \item 综合算式： $50 - 4 \times 6 = 50 - 24 = 26$ （元）。
\end{enumerate}


\end{answer}
